\documentclass{article}

\usepackage{amsmath}
\usepackage{amssymb}
\usepackage{xcolor}
\usepackage{tikz}
\usepackage[shortlabels]{enumitem}

\usepackage[utf8]{inputenc}
\usepackage[T1]{fontenc}

\usepackage[polish]{babel}

\usepackage{listings}
\lstset{
    language=SQL,
    inputencoding=utf8x,
    extendedchars=\true,
    literate={ą}{{\k{a}}}1
             {Ą}{{\k{A}}}1
             {ę}{{\k{e}}}1
             {Ę}{{\k{E}}}1
             {ó}{{\'o}}1
             {Ó}{{\'O}}1
             {ś}{{\'s}}1
             {Ś}{{\'S}}1
             {ł}{{\l{}}}1
             {Ł}{{\L{}}}1
             {ż}{{\.z}}1
             {Ż}{{\.Z}}1
             {ź}{{\'z}}1
             {Ź}{{\'Z}}1
             {ć}{{\'c}}1
             {Ć}{{\'C}}1
             {ń}{{\'n}}1
             {Ń}{{\'N}}1
}

\title{Przestrzenie wektorowe - zestaw III}
\date{2026}

\begin{document}
\maketitle

\begin{enumerate}

\item 
Wyznaczyć postać kanoniczną Jordana macierzy:
        \begin{enumerate}[a)]
            \item $A = \begin{bmatrix}\begin{array}{rrr} 1 & 2 & 0\\ 2 & 1 & 0\\ 0 & 0 & 2 \end{array}\end{bmatrix}$,
            \item $A = \begin{bmatrix}\begin{array}{rrr} 3 & 0 & -1\\ 0 & 3 & -1\\ 0 & 1 & 1 \end{array}\end{bmatrix}$,
            \item $A = \begin{bmatrix}\begin{array}{rrr} -1 & 1 & 1\\ 1 & -1 & 1\\ 1 & 1 & -1 \end{array}\end{bmatrix}$,
            \item $A = \begin{bmatrix}\begin{array}{rrrr} 2 & 0 & 0 & 0\\ 0 & 1 & 1 & 0\\ 0 & 0 & 1 & 0\\ 0 & 2 & -1 & 2 \end{array}\end{bmatrix}$.
        \end{enumerate}
        Do wyznaczenia wielomianu charakterystycznego macierzy $A$: $|\lambda I - A| = \lambda^{n} + a_{n-1}\lambda^{n-1} + \cdots + a_{0}$ można wykorzystać zależność:
        $$
            a_{k} = (-1)^{n-k}\sum_{j=1}^{ {n}\choose{k}} {M_{j}(n-k)},\,k=0,1,\cdots,(n-1)
        $$
        gdzie $M_{j}(n-k)$ jest $j$-tym minorem głównym rzędu $n-k$ macierzy $A$, tj. wyznacznikiem macierzy powstającej z $A$ przez wykreślenie $k$ kolumn i wierszy o tych samych numerach.
        
        Na przyjkład dla macierzy $A = \begin{bmatrix}\begin{array}{rrr} 2 & 1 & 2\\ 1 & 3 & 1\\ 1 & 1 & 2 \end{array}\end{bmatrix}$ mamy:
        $|\lambda I - A| = \begin{vmatrix}\begin{array}{rrr} \lambda - 2 & -1 & -2\\ -1 & \lambda - 3 & -1\\ -1 & -1 & \lambda - 2 \end{array}\end{vmatrix} = \lambda^{3} + a_{2}\lambda^{2} + a_{1}\lambda + a_{0}$, gdzie:
        \begin{itemize}
            \item $a_{2} = (-1)^{3-2} \sum_{j=1}^{{{3}\choose{2}}=3}{M_{j}(3-2)} = -(2 + 3 + 2) = -7$,
            \item $a_{2} = (-1)^{3-1} \sum_{j=1}^{{{3}\choose{1}}=3}{M_{j}(3-1)} = \begin{vmatrix}\begin{array}{rr} 2 & 1\\ 1 & 3 \end{array}\end{vmatrix} + \begin{vmatrix}\begin{array}{rr} 2 & 2\\ 1 & 2 \end{array}\end{vmatrix} + \begin{vmatrix}\begin{array}{rr} 3 & 1\\ 1 & 2 \end{array}\end{vmatrix} = 12$,
            \item $a_{2} = (-1)^{3-0} \sum_{j=1}^{{{3}\choose{0}}=1}{M_{j}(3-0)} = \begin{vmatrix}\begin{array}{rrr} 2 & 1 & 2\\ 1 & 3 & 1\\ 1 & 1 & 2 \end{array}\end{vmatrix} = -5$,
        \end{itemize}

\item 
Wyznaczyć macierz $e^{At}$, gdy:
        \begin{enumerate}[a)]
            \item $A = \begin{bmatrix}\begin{array}{rrr} \lambda & 1 & 0\\ 0 & \lambda & 1\\ 0 & 0 & \lambda \end{array}\end{bmatrix}$,
            \item $A = \begin{bmatrix}\begin{array}{rrr} \lambda_{1} & 1 & 0\\ 0 & \lambda_{1} & 0\\ 0 & 0 & \lambda_{2} \end{array}\end{bmatrix}$,
            \item $A = \begin{bmatrix}\begin{array}{rr} 1 & 2\\ 2 & 1 \end{array}\end{bmatrix}$,
            \item $A = \begin{bmatrix}\begin{array}{rrr} 3 & 0 & 0\\ -2 & 1 & 0\\ 1 & 1 & 1 \end{array}\end{bmatrix}$.
        \end{enumerate}

\item 
Wyznaczyć macierz $A^{-1}$ wykorzystując twierdzenie Cayley'a-Hamiltona (czyli fakt, że macierz $A$ spełnia swoje równanie charakterystyczne):
        \begin{enumerate}[a)]
            \item $A = \begin{bmatrix}\begin{array}{rrr} 2 & 1 & 2\\ 1 & 3 & 1\\ 1 & 1 & 2 \end{array}\end{bmatrix}$,
            \item $A = \begin{bmatrix}\begin{array}{rrr} 2 & 2 & 4\\ 1 & 2 & 2\\ 2 & 1 & 2 \end{array}\end{bmatrix}$,
            \item $A = \begin{bmatrix}\begin{array}{rrrr} 2 & 2 & 4 & 2\\ 1 & 2 & 2 & 0\\ 2 & 1 & 2 & 1 \\ 1 & 2 & 3 & 4\end{array}\end{bmatrix}$,
            \item $A = \begin{bmatrix}\begin{array}{rrrr} 1 & 0 & 3 & 0\\ 0 & 1 & 0 & 3\\ 3 & 0 & 1 & 0 \\ 0 & 3 & 0 & 1\end{array}\end{bmatrix}$,
            %\item $A = \begin{bmatrix}\begin{array}{rrrrr} 1 & 2 & 3 & 4 & 5\\ 5 & 1 & 2 & 3 & 4\\ 4 & 5 & 1 & 2 & 3 \\ 3 & 4 & 5 & 1 & 2 \\ 2 & 3 & 4 & 5 & 1\end{array}\end{bmatrix}$.
        \end{enumerate} 

\item 
Wykazać, że zachodzi własność:
    \begin{enumerate}[a)]
        \item $e^{A^{T}} = \left( e^{A}\right)^{T} $,
        \item jeżeli $A=A^{T}$, to $e^{A}$ jest symetryczna.
    \end{enumerate}
\end{enumerate}


\newpage
\textbf{Odpowiedzi:}
\begin{description}
    \item [1:] a) $\begin{bmatrix}\begin{array}{rrr} -1 & 0 & 0\\ 0 & 2 & 0\\ 0 & 0 & 3 \end{array}\end{bmatrix}$ (macierz modalna: $T = \begin{bmatrix}\begin{array}{rrr} 1 & 0 & 1\\ -1 & 0 & 1\\ 0 & 1 & 0 \end{array}\end{bmatrix}$), b) $\begin{bmatrix}\begin{array}{rrr} 2 & 1 & 0\\ 0 & 2 & 0\\ 0 & 0 & 3 \end{array}\end{bmatrix}$  (macierz modalna: $T = \begin{bmatrix}\begin{array}{rrr} 1 & 1 & 1\\ 1 & 1 & 0\\ 1 & 0 & 0 \end{array}\end{bmatrix}$), c) $\begin{bmatrix}\begin{array}{rrr} -2 & 1 & 0\\ 0 & -2 & 0\\ 0 & 0 & 1 \end{array}\end{bmatrix}$ (macierz modalna: $T = \begin{bmatrix}\begin{array}{rrr} -1 & -1 & 1\\ 1 & 0 & 1\\ 0 & 1 & 1 \end{array}\end{bmatrix}$), d) $\begin{bmatrix}\begin{array}{rrrr} 1 & 1 & 0 & 0\\ 0 & 1 & 0 & 0\\ 0 & 0 & 2 & 0\\ 0 & 0 & 0 & 2 \end{array}\end{bmatrix}$ (macierz modalna $T = \begin{bmatrix}\begin{array}{rrrr} 0 & 0 & 1 & 0\\ 1 & 0 & 0 & 0\\ 0 & 1 & 0 & 0\\ -2 & -1 & 0 & 1 \end{array}\end{bmatrix}$)

    \item [2:] a) $e^{At} = \begin{bmatrix}\begin{array}{rrr} e^{\lambda t} & t e^{\lambda t} & \frac{1}{2}t^{2} e^{\lambda t}\\ 0 & e^{\lambda t} & t e^{\lambda t}\\ 0 & 0 & e^{\lambda t} \end{array}\end{bmatrix}$,
    b) $e^{At} = \begin{bmatrix}\begin{array}{rrr} e^{\lambda_{1} t} & t e^{\lambda_{1} t} & 0\\ 0 & e^{\lambda_{1} t} & 0\\ 0 & 0 & e^{\lambda_{2} t} \end{array}\end{bmatrix}$,
    c)  $e^{At} = \frac{1}{2} \begin{bmatrix}\begin{array}{rr} e^{3t} + e^{-t} & e^{3t} - e^{-t}\\ e^{3t} - e^{-t} & e^{3t} + e^{-t} \end{array}\end{bmatrix}$,
    d)  $\begin{bmatrix}\begin{array}{rrr} {\mathrm{e}}^{3\,t} & 0 & 0\\ {\mathrm{e}}^t-{\mathrm{e}}^{3\,t} & {\mathrm{e}}^t & 0\\ t\,{\mathrm{e}}^t & t\,{\mathrm{e}}^t & {\mathrm{e}}^t \end{array}\end{bmatrix}$

    \item [3:] a) $A^{-1} = \begin{bmatrix}\begin{array}{rrr} 1 & 0 & -1\\ -\frac{1}{5} & \frac{2}{5} & 0\\ -\frac{2}{5} & -\frac{1}{5} & 1 \end{array}\end{bmatrix}$,
    b) $A^{-1} = \begin{bmatrix}\begin{array}{rrr} -\frac{1}{2} & 0 & 1\\ -\frac{1}{2} & -1 & 0\\ -\frac{3}{4} & -\frac{1}{2} & -\frac{1}{2} \end{array}\end{bmatrix}$
    c) $A^{-1} = \begin{bmatrix}\begin{array}{rrrr} -\frac{1}{2} & 0 & 1 & 0\\ -\frac{3}{4} & \frac{5}{6} & \frac{1}{6} & \frac{1}{3}\\ 1 & -\frac{1}{3} & -\frac{2}{3} & -\frac{1}{3}\\ -\frac{1}{4} & -\frac{1}{6} & \frac{1}{6} & \frac{1}{3} \end{array}\end{bmatrix}$, potęgi $A$: $A^{2} = \begin{bmatrix}\begin{array}{rrrr} 16 & 16 & 26 & 16\\ 8 & 8 & 12 & 4\\ 10 & 10 & 17 & 10\\ 14 & 17 & 26 & 21 \end{array}\end{bmatrix}$, $A^{3} = \begin{bmatrix}\begin{array}{rrrr} 116 & 122 & 196 & 122\\ 52 & 52 & 84 & 44\\ 74 & 77 & 124 & 77\\ 118 & 130 & 205 & 138 \end{array}\end{bmatrix}$,
    d) $A^{-1} = \begin{bmatrix}\begin{array}{rrrr} -\frac{1}{8} & 0 & \frac{3}{8} & 0\\ 0 & -\frac{1}{8} & 0 & \frac{3}{8}\\ \frac{3}{8} & 0 & -\frac{1}{8} & 0\\ 0 & \frac{3}{8} & 0 & -\frac{1}{8} \end{array}\end{bmatrix}$, potęgi $A$: $A^{2} = \begin{bmatrix}\begin{array}{rrrr} 10 & 0 & 6 & 0\\ 0 & 10 & 0 & 6\\ 6 & 0 & 10 & 0\\ 0 & 6 & 0 & 10 \end{array}\end{bmatrix}$, $A^{3} = \begin{bmatrix}\begin{array}{rrrr} 28 & 0 & 36 & 0\\ 0 & 28 & 0 & 36\\ 36 & 0 & 28 & 0\\ 0 & 36 & 0 & 28 \end{array}\end{bmatrix}$
\end{description}

\end{document}